% Standalone solution dossier for prop:a5-partial-funnel.
\documentclass[../main.tex]{subfiles}
\begin{document}

% === SOLUTION HEADER (proof provenance) =====================================
% ledger-node: prop:a5-partial-funnel
% refines: prop:a5-partial-funnel
% depends_on: lem:a5-pi-exp-tail
% bounded_by: obs:heavy-tail-no-classical
% evidence_run: none
% checked_by: agent
% author: /root/a4_a5
% proof_completion: /root/review_a4_a5
% dossier_editor: /root/review_a4_a5
% reviewer: /root/cross_review
% review: research/reviews/2026-08-21-a5-funnel-second-agent-audit.md
% date: 2026-08-21
% ============================================================================

\section*{Solution: the partial Neal-funnel phase}

\begin{lemma}
If a law satisfies Poincar\'e with finite constant $C$, every real $1$-Lipschitz $f\in L^2$
has a finite exponential moment about a median for all sufficiently small positive exponents.
\end{lemma}

\begin{proof}
Let $m$ be a median and $a=2\sqrt C$. For $t\ge0$, apply Poincar\'e to
$g=\min\{(f-m-t)_+,a\}$. The independent-copy variance identity and
$\P(f\le m+t)\ge1/2$ give
\[
 \tfrac12a^2\P(f\ge m+t+a)\le\Var(g)
 \le C\P(m+t<f<m+t+a)\le C\P(f>m+t).
\]
Thus each increase by $a$ halves the upper tail. Applying the same argument to $-f$ yields a
two-sided exponential tail and hence a finite exponential moment for
$c<(\log2)/a$.
\end{proof}

\begin{proposition}
For $s>0$, let $u\sim N(0,s^2)$ and $z\sim N(0,1)$ be independent, let $\theta=e^uz$, and define
$z_\alpha=e^{-\alpha u}\theta=e^{(1-\alpha)u}z$ for $0\le\alpha\le1$. In the Euclidean
coordinates $(u,z_\alpha)$,
\[
 C_P=+\infty\quad(0\le\alpha<1),\qquad
 C_P=\max\{s^2,1\}\quad(\alpha=1).
\]
\end{proposition}

\begin{proof}
Fix $\alpha<1$ and put $\beta=1-\alpha>0$. The coordinate $z_\alpha$ is one-Lipschitz in the
Euclidean $(u,z_\alpha)$ space and lies in $L^2$, since
$\E z_\alpha^2=e^{2\beta^2s^2}$. For every $c>0$,
\[
\E e^{c|z_\alpha|}
 \ge\P(|z|\ge1)\E[e^{ce^{\beta u}}\mathbf1_{\{u>0\}}]=\infty.
\]
The last integral diverges because
$ce^{\beta u}-u^2/(2s^2)\to+\infty$. An exponential moment about a median would imply an
uncentred one by the triangle inequality, so the lemma contradicts any finite Poincar\'e constant.
At $\alpha=1$, $z_\alpha=z$ and the law is
the product $N(0,s^2)\otimes N(0,1)$, whose Poincar\'e constant is
$\max\{s^2,1\}$ by tensorization.
\end{proof}

\paragraph{Audit state.}
The completed concentration and funnel proof received a second independent agent audit after the
$s>0$ hypothesis and all integrability steps were made explicit.

\end{document}
